\section{问题三：全域达标条件与烘干时间}
\subsection{从空间约束到事件函数}
沿用问题二的固定域耦合方程，前4 h使用观测环境，后期使用式\eqref{eq_late}。为表达题目“各处低于0.15 kg/kg”的要求，定义
\begin{equation}\label{eq_event}
M(t)=\max_{0\le r\le R}C(r,t),\qquad
g(t)=M(t)-0.15,\qquad
t_* =\inf\{t\ge0:M(t)<0.15\}.
\end{equation}
在连续向下穿越的情况下，式\eqref{eq_event}定义的临界时刻满足$M(t_*)=0.15$。严格低于阈值的时间集合通常没有最小元素，因此另给稍晚的报告时刻$t_{\mathrm{rep}}$，并直接复核未舍入的$M(t_{\mathrm{rep}})<0.15$。

计算中使用全部内部节点的最大值$M_h(t)=\max_iC_i(t)$，记录最大值位置并检查径向分布。只有数值剖面支持中心最湿时，才用中心曲线解释控制位置；不事先用中心值替代全域搜索。题目要求的5个表格位置和21个Excel位置均是输出采样，不承担全域最大值检查。

\subsection{连续事件定位与严格裕量}
采用$G_8$网格联合积分；相对容差$10^{-9}$，温度、含水率绝对容差分别为$10^{-8}$、$10^{-10}$。前4 h最大步长2 s，之后60 s，在输入切换处分段接续。搜索$g$由正到负的括号后，在事件附近以最大0.5 s步长局部精算，再用带括号求根定位临界时刻。

为避免小幅离散误差改变严格不等式，在同一物理时刻$t_{*,h}$比较生产解、细网格解和时间收紧解。令
\begin{equation}\label{eq_margin}
\varepsilon_C=2\max\left\{|M_h(t_{*,h})-M_{h/2}(t_{*,h})|,
|M_h(t_{*,h})-M_{h,\mathrm{tight}}(t_{*,h})|\right\}+10^{-8}.
\end{equation}
本问式\eqref{eq_margin}给出$\varepsilon_C=3.5088\times10^{-7}$ kg/kg。先寻找$M_h<0.15-\varepsilon_C$的时刻，再向上选择整秒并向上保留四位小时，最终在三组解上复核。该裕量反映已完成的加密比较，是经验数值裕量，不是连续PDE解的严格误差界。

生产网格临界时刻为57.4725195022 h，最终报告
\begin{equation}\label{eq_q3_answer}
\boxed{t_{\mathrm{rep},3}=57.4731\ \mathrm h=206903.16\ \mathrm s.}
\end{equation}
式\eqref{eq_q3_answer}对应的生产解最大含水率为0.1499993913 kg/kg；细网格和时间收紧解在同一时刻的最大值分别为0.1499992209、0.1499993913 kg/kg，均严格低于0.15。最大值位于中心，保存剖面从中心向外不增。

\subsection{径向分布与后期干燥特征}
表\ref{tab_q3C}给出每6 h及最终报告时刻的水分浓度，图\ref{fig_q3history}展示最大含水率随时间的变化，图\ref{fig_q3profiles}展示径向剖面。
\input{tables/q3C}
\figone{q3_C_max_history.pdf}{.78}{全域最大含水率逐步接近干燥阈值。后期下降速度减慢，表面先达标不能代表内部已完成干燥。}{fig_q3history}
\figone{q3_C_profiles.pdf}{.78}{固定半径模型每6 h及终时的含水率剖面。中心保持较高含水率，最终由中心附近控制全域达标。}{fig_q3profiles}

例如，18 h时表面已降至0.0845 kg/kg，而中心仍为0.2988 kg/kg；48 h时中心仍为0.1618 kg/kg，尚未达标。这表明单凭表面测量或平均含水率判定结束会遗漏内部湿区。后期温度接近恒定后，失水使局部扩散系数继续下降，因而不能把早期失水速率线性外推到阈值。

表\ref{tab_q3C}末行中心显示0.1500，是0.1499993913保留四位小数的结果；达标判断使用未舍入值，不以显示值是否为0.1499为准。完整结果文件从60 s起每60 s保存一次，另追加206903.16 s的准确终时，不将最近整分钟冒充最终报告时刻。
\FloatBarrier
